% Chapter 12 of Optimizers in the Wild: learning-rate schedules.
\newcommand{\schedulesG}{\mathcal{G}}
\newcommand{\schedulesK}{\mathcal{K}}

\casechapter{Schedules}{Cosine, warmup--stable--decay and the river valley}{Cosine and WSD schedules}{ch:schedules}
\begin{refsection}[chapters/12-schedules.bib]

\begin{chaptersummary}
Large models are trained with a learning rate that rises briefly and then decays, along a cosine, along a straight line, or in a short cooldown after a long plateau, and the loss falls fastest exactly when the learning rate falls. For convex Lipschitz losses we prove a bound on the last iterate of stochastic gradient descent under an arbitrary schedule, show that linear decay to zero attains the rate of iterate averaging exactly, and evaluate the bound for the schedules in use. In a river valley, a slowly descending direction flanked by stiff and noisy ones, we compute the expected loss exactly for every schedule: it is a term that depends on the area under the schedule plus a term that remembers only the recent schedule, which is the form of the empirical laws that predict loss curves across schedules, and it predicts how the best cooldown scales with the length of the run. Neither idealization has a landscape that changes during training, and the river valley's predictions depend on constants that must be measured.
\end{chaptersummary}

\section{In the wild: the shape of the learning rate}\label{sec:schedules:wild}

Ask what learning rate a large language model was trained with and the answer is a curve. Most published pretraining recipes since GPT-3 warm the learning rate up over a small part of the run and then decay it to a small value at the planned end, and the shape of the decay is one of the few hyperparameters still argued about in public. As of 2025 three shapes dominate the open recipes.

\begin{practice}[GPT-3, \textcite{BrownEtAl2020}; LLaMA, \textcite{TouvronEtAl2023}; Chinchilla, \textcite{HoffmannEtAl2022}; MiniCPM, \textcite{HuEtAl2024}; \textcite{HageleEtAl2024}; \textcite{BergsmaEtAl2025}]
Warm the learning rate up linearly from zero to its peak $\eta_{\max}$, over the first $375$ million tokens (GPT-3) or $2000$ steps (LLaMA). Then do one of three things. \emph{Cosine}: decay along half a period of a cosine to a tenth of the peak, $\eta_t=\eta_{\max}\bigl(0.1+0.45(1+\cos(\pi t/T))\bigr)$; LLaMA and Chinchilla reach the tenth at the end of the run, GPT-3 after $260$ of its $300$ billion tokens. \emph{Linear decay to zero}, $\eta_t=\eta_{\max}(1-t/T)$ \parencite{DefazioEtAl2023,BergsmaEtAl2025}. \emph{Warmup--stable--decay} (WSD): hold the peak, then decay to zero or nearly zero over the last $10$--$20\%$ of the steps (MiniCPM; \textcite{HageleEtAl2024} recommend about $20\%$, with the shape $1-\sqrt{\cdot}$). The original transformer used the inverse square root, $\eta_t\propto\min(t^{-1/2},\,t\,T_{\mathrm w}^{-3/2})$ with $T_{\mathrm w}=4000$ \parencite{VaswaniEtAl2017}.
\end{practice}

The cosine\index{cosine schedule} came from vision, where \textcite{LoshchilovHutter2017} introduced cosine annealing with periodic restarts; language models kept a single half period.\sidenote{Earlier vision recipes divided the learning rate by $10$ at two or three fixed epochs, a \emph{step decay}. The classical rule of stochastic approximation is $\eta_t\propto1/t$, which satisfies the conditions $\sum_t\eta_t=\infty$ and $\sum_t\eta_t^2<\infty$ of \textcite{RobbinsMonro1951}. Every schedule in the box violates the second: at a fixed peak, $\sum_{t\le T}\eta_t^2$ grows without bound with the horizon $T$, linearly for the decaying shapes and like $\log T$ for the inverse square root.} The shapes look very different, yet their final losses in tuned runs are close, and which is best is still debated.

\begin{observation}[The cooldown drop; \textcite{HuEtAl2024}, \textcite{HageleEtAl2024}]\label[observation]{obs:schedules:cooldown}
Train with WSD\index{warmup--stable--decay (WSD)}\index{cooldown}. During the stable phase the loss decreases slowly; when the decay begins it falls much faster, and after a decay over $10$--$20\%$ of the steps it matches, or slightly beats, a cosine schedule of the same length. Longer decays help a little more, with diminishing returns.
\end{observation}

\begin{observation}[The horizon; \textcite{HoffmannEtAl2022}]\label[observation]{obs:schedules:horizon}
A cosine schedule planned for more steps than the run takes, by more than about $25\%$, gives a worse model than one planned for exactly the steps taken. WSD runs of different lengths coincide until their cooldowns, so one stable phase can be branched into cooldowns of several lengths, which makes scaling experiments cheaper \parencite{HageleEtAl2024}.
\end{observation}

The next two observations are of a different kind: fits and analogies that describe whole loss curves.

\begin{observation}[Loss curves across schedules; \textcite{TissueEtAl2025}, \textcite{LuoEtAl2025}]\label[observation]{obs:schedules:laws}
Fitted to the loss curves of one or two schedules, a law of the form
\[
  \Loss(t)\approx L_\infty+C\,S_1(t)^{-\alpha}-\mathrm{LD}(t),\qquad S_1(t)=\sum_{s\le t}\eta_s ,
\]
predicts the loss curves of the same model under other schedules\index{loss-curve laws}. The \emph{forward area} $S_1(t)$ is the area under the schedule so far; the \emph{loss drop} $\mathrm{LD}(t)$ adds, for each decrease of the learning rate, a term proportional to the decrease that grows and then saturates with the number of steps \parencite{TissueEtAl2025} or with the area \parencite{LuoEtAl2025} since the decrease.
\end{observation}

\begin{observation}[Convex theory agrees; \textcite{SchaippEtAl2025}]\label[observation]{obs:schedules:convex}
Evaluate, for each schedule and each step $t$, the bound on $\E\Loss(\theta_t)-\Loss^\star$ that convex optimization gives for stochastic gradient descent on a convex Lipschitz loss. The curves look like language-model loss curves: WSD plateaus and then drops suddenly, and cosine and WSD end close together. \Textcite{SchaippEtAl2025} report that the bound reproduces the matching performance of cosine and WSD and the necessity and length of the cooldown, and that learning rates chosen with it improved Llama-type models with $124$M and $210$M parameters.
\end{observation}

Language-model training is far from convex, so \cref{obs:schedules:convex} suggests that the schedule acts through something simpler than the landscape; in the experiments of \textcite{DefazioEtAl2023}, too, warmup followed by linear decay matched or beat cosine almost everywhere, as their worst-case theory predicts. Why does the loss fall when the learning rate falls, why do such different shapes end so close together, and how long must the decay be?

\section{The idealization: two pictures of a schedule}\label{sec:schedules:ideal}

Throughout, $\theta_1$ is the initial point and stochastic gradient descent with step sizes $\eta_1,\eta_2,\dots$ is
\begin{equation}\label{eq:schedules:sgd}
  \theta_{t+1}=\theta_t-\eta_t\,g_t ,\qquad t=1,2,\dots,
\end{equation}
with $g_t$ a stochastic gradient at $\theta_t$; ``learning rate'' (in practice) and ``step size'' (in the theory) are synonyms. Let $\cF_t$ be the $\sigma$-algebra generated by $g_1,\dots,g_{t-1}$, so that $\theta_t$ is $\cF_t$-measurable. A \term{schedule}\index{schedule} is the sequence $(\eta_t)$, often written $\eta_t=\eta\,\varphi(t/T)$ with a peak $\eta$ and a \emph{shape} $\varphi\colon[0,1]\to[0,1]$. The \term{area}\index{area under the schedule} before step $t$ is $A_t=\sum_{s=1}^{t-1}\eta_s$, the time of the gradient flow that the run follows when its steps are small.

\paragraph{Convex Lipschitz losses.} The first picture keeps nothing of the landscape except convexity.

\begin{assumption}[Convex Lipschitz stochastic optimization]\label[assumption]{ass:schedules:convex}
The loss $\Loss\colon\R^d\to\R$ is convex and has a minimizer $\theta^\star$, and $D=\norm{\theta_1-\theta^\star}_2$. The stochastic gradients satisfy $\E[g_t\mid\cF_t]\in\partial\Loss(\theta_t)$ and $\E\bigl[\norm{g_t}_2^2\mid\cF_t\bigr]\le G^2$.
\end{assumption}

The constant $G$ bounds subgradients and noise together, so the subgradient method on a $G$-Lipschitz function and stochastic gradient descent on a smooth function with $\norm{\nabla\Loss}_2^2+\E\norm{\zeta_t}_2^2\le G^2$ are both covered; the norm is Euclidean because stochastic gradient descent is steepest descent in $\ell_2$ (\cref{ch:steepest}). A statement about the whole class cannot see the scales of a particular landscape, and this will show in the numbers.

\paragraph{The river valley.} The second picture keeps the simplest landscape with two time scales. \Textcite{WenEtAl2025} proposed that the pretraining loss of a language model is shaped like a \term{river valley}\index{river valley}: a deep valley whose floor descends slowly along a river, between steep walls. With a large learning rate the iterate oscillates between the walls while it moves quickly along the river; when the learning rate decays the oscillation dies out, and the progress made along the river shows in the loss.\sidenote{\Textcite{WenEtAl2025} derive the picture in a toy model trained on bigram data and use it to design a WSD variant that reuses decayed checkpoints. The quadratic walls and the noise below are our simplification.} We make the walls quadratic, the oscillation noise-driven, and the river straight or weakly curved.

\begin{assumption}[The river valley]\label[assumption]{ass:schedules:river}
The parameter is $\theta=(x,y)$ with river coordinates $x\in\R^p$ and hill coordinates $y\in\R^m$, and
\[
  \Loss(x,y)=r(x)+\frac12\sum_{i=1}^m h_i\,y_i^2,\qquad h_i>0 .
\]
The stochastic gradient is $g_t=\nabla\Loss(\theta_t)+\zeta_t$ with $\E[\zeta_t\mid\cF_t]=0$ and $\E[\zeta_{t,i}^2\mid\cF_t]=\sigma_i^2$ for each hill coordinate $y_i$. The river $r$ is linear, with a constant gradient $\nabla r$, or a quadratic whose curvatures are much smaller than the $h_i$, with noise of fixed variance in each of its eigendirections.
\end{assumption}

With a quadratic river the model is the noisy quadratic model\index{noisy quadratic model} of \cref{ch:noise}, its spectrum split into hills, whose curvature times the area of the run is large, and river directions, for which it is small. In a hill the iterate soon forgets its past and sits in a noise-driven equilibrium; along the river it never forgets, and it is still descending when the run ends.

\begin{dictionary}
step size $\eta_t$, peak $\eta$, shape $\varphi$ & the schedule $\eta_{\max}\varphi(t/T)$\\
area $A_t=\sum_{s<t}\eta_s$ & the forward area $S_1$ of the empirical laws; how far the run has travelled along the river\\
river $r(x)$ & the part of the loss that keeps falling at a constant learning rate, roughly a power of the area\\
hill $y_i$, curvature $h_i$ & a sharp direction of the Hessian, along which the iterate oscillates at a high learning rate\\
noise variance $\sigma_i^2$ & gradient noise along that direction, $\Sigma_{ii}/B$ at batch size $B$\\
noise floor $\eta\sigma_i^2/(2(2-\eta h_i))$ & the excess loss of the stable phase, removed by the cooldown\\
relaxation time $1/(2\eta h_i)$ steps & how long the loss takes to respond to a change of the learning rate\\
annealing term $\sum_k(\eta_{k-1}-\eta_k)\,\schedulesG(A_T-A_k)$ & the loss drop $\mathrm{LD}(t)$ of the empirical laws\\
$D$ and $G^2$ & the distance to a good solution; the mean squared gradient
\end{dictionary}

\paragraph{What was thrown away.} The convex picture throws away curvature, and with it every distinction between fast and slow directions, as well as nonconvexity and every optimizer but plain stochastic gradient descent; it comes first because in it a statement about every schedule and every loss of a class can be proved. The river valley throws away the coupling of river and hills (in the picture of \textcite{WenEtAl2025} the river bends and the walls turn with it), changes of curvature, noise that depends on the position, and momentum, Adam and weight decay. It keeps what the convex picture lacks, two time scales, and it is solvable exactly.

\paragraph{What one had to notice.} The loss at the end of a run is a functional of the whole schedule, and in both pictures it depends on the schedule through two features with very different memories: the area, which remembers everything, and an average of the step sizes over the last few relaxation times, which remembers almost nothing. Large steps buy area; small steps at the end erase the short memory. A schedule is a compromise between two terms, and the shapes in use are different compromises between the same two terms.

\section{The theory: the last iterate under an arbitrary schedule}\label{sec:schedules:convex}

The classical analysis bounds the average iterate: by the online-gradient-descent bound of \cref{thm:toolkit:ogd} and Jensen's inequality, the $\eta$-weighted average of the iterates is within $(D^2+G^2\sum_t\eta_t^2)/(2\sum_t\eta_t)$ of the optimum, which is $DG/\sqrt T$ for the constant step $D/(G\sqrt T)$ \parencite{NemirovskiEtAl2009}. Practitioners use the \term{last iterate}\index{last iterate}.\sidenote{Averaging in stochastic approximation goes back to \textcite{PolyakJuditsky1992}; in deep learning it survives as the moving average of the weights and in schedule-free methods (\cref{ch:momentum}).} The tool for $\theta_T$ is the same expansion with the minimizer replaced by an earlier iterate.\sidenote{The proofs use one more stochastic gradient $g_T$ at $\theta_T$ and the point $\theta_{T+1}$. Neither affects $\theta_1,\dots,\theta_T$, so the step size $\eta_T$ is a free parameter of the bounds.}

\begin{lemma}[Anchored regret]\label[lemma]{lem:schedules:anchor}
Let \cref{ass:schedules:convex} hold and let $\eta_1,\dots,\eta_T>0$. For each $k\in\{1,\dots,T\}$ and for $u=\theta^\star$ or $u=\theta_k$,
\[
  \sum_{t=k}^{T}\eta_t\,\E\bigl[\Loss(\theta_t)-\Loss(u)\bigr]\le\frac12\,\E\norm{\theta_k-u}_2^2+\frac{G^2}2\sum_{t=k}^{T}\eta_t^2 .
\]
\end{lemma}

\begin{proof}
The plan is to expand the distance of $\theta_{t+1}$ to $u$ and telescope from $t=k$ on. Let $s_t=\E[g_t\mid\cF_t]\in\partial\Loss(\theta_t)$; by Jensen's inequality $\norm{s_t}_2\le G$. All expectations are finite: $(\E\norm{\theta_t}_2^2)^{1/2}\le\norm{\theta_1}_2+G\sum_{s<t}\eta_s$ by Minkowski's inequality, and $0\le\Loss(\theta_t)-\Loss^\star\le\Loss(\theta_1)-\Loss^\star+G\norm{\theta_t-\theta_1}_2$ by the subgradient inequality at $\theta_t$. By \eqref{eq:schedules:sgd},
\[
  \norm{\theta_{t+1}-u}_2^2=\norm{\theta_t-u}_2^2-2\eta_t\ip{g_t,\theta_t-u}+\eta_t^2\norm{g_t}_2^2 .
\]
For $t\ge k$ both $\theta_t$ and $u$ are $\cF_t$-measurable, so conditioning on $\cF_t$ replaces $g_t$ by $s_t$ in the middle term, and $\ip{s_t,\theta_t-u}\ge\Loss(\theta_t)-\Loss(u)$. Hence $2\eta_t\E[\Loss(\theta_t)-\Loss(u)]\le\E\norm{\theta_t-u}_2^2-\E\norm{\theta_{t+1}-u}_2^2+\eta_t^2G^2$. Summing over $t=k,\dots,T$ and dropping $-\E\norm{\theta_{T+1}-u}_2^2\le0$ gives the claim.
\end{proof}

\begin{theorem}[Last iterate under an arbitrary schedule]\label{thm:schedules:lastiterate}
Let \cref{ass:schedules:convex} hold, let $\eta_1,\dots,\eta_T>0$, and let $\theta_1,\dots,\theta_T$ be generated by \eqref{eq:schedules:sgd}. Write $R_k=\sum_{t=k}^{T}\eta_t$ and $V_k=\sum_{t=k}^{T}\eta_t^2$. Then
\begin{equation}\label{eq:schedules:lastiterate}
  \E\Loss(\theta_T)-\Loss^\star\ \le\ \frac{D^2+G^2V_1}{2R_1}\;+\;\frac{G^2}{2}\sum_{k=1}^{T-1}\frac{\eta_k}{R_{k+1}}\cdot\frac{V_k}{R_k} .
\end{equation}
\end{theorem}

\begin{proof}
The plan is to compare, for each $k$, the $\eta$-weighted mean of the expected losses from step $k$ on with the loss at step $k$, using $\theta_k$ as the comparator in \cref{lem:schedules:anchor}, and to telescope from the mean over the whole run down to the last iterate. Write $\ell_t=\E\Loss(\theta_t)$ and $\bar\ell_k=R_k^{-1}\sum_{t=k}^{T}\eta_t\ell_t$, so that $\bar\ell_T=\ell_T$.

\emph{The whole run.} \Cref{lem:schedules:anchor} with $k=1$ and $u=\theta^\star$ gives $R_1(\bar\ell_1-\Loss^\star)\le\frac12(D^2+G^2V_1)$.

\emph{One step of the telescope.} Let $1\le k\le T-1$. \Cref{lem:schedules:anchor} with $u=\theta_k$ gives $R_k(\bar\ell_k-\ell_k)\le\frac12G^2V_k$. By definition $R_k\bar\ell_k=\eta_k\ell_k+R_{k+1}\bar\ell_{k+1}$; subtracting $R_{k+1}\bar\ell_k=(R_k-\eta_k)\bar\ell_k$ from both sides,
\[
  R_{k+1}\bigl(\bar\ell_{k+1}-\bar\ell_k\bigr)=\eta_k\bigl(\bar\ell_k-\ell_k\bigr)\le\frac{G^2}2\,\frac{\eta_kV_k}{R_k} .
\]
\emph{Summation.} Divide by $R_{k+1}>0$ and add the identity $\ell_T-\Loss^\star=(\bar\ell_1-\Loss^\star)+\sum_{k=1}^{T-1}(\bar\ell_{k+1}-\bar\ell_k)$.
\end{proof}

The first term of \eqref{eq:schedules:lastiterate} is the averaging bound. The second, the price of the last iterate, averages the step sizes near the end: $V_k/R_k$ is the mean step size from step $k$ on, and its weight $\eta_k/R_{k+1}$, the current step over the area still to come, is negligible early and of order one at the end. The two terms are the convex counterparts of the river and the hills of \cref{sec:schedules:river}: one decreases with the area, the other remembers the recent schedule. Constant steps make the weights $1/(T-k)$, whose harmonic sum costs a logarithm; a decay to zero makes the step sizes small where the weights are large.\sidenote{The comparison with an earlier iterate is the argument of \textcite{ShamirZhang2013}; \textcite{Orabona2020} gave a version without a bounded domain, which we follow. \Textcite{DefazioEtAl2023} prove a last-iterate bound for general schedules by a reduction from regret, and \textcite{SchaippEtAl2025} evaluate such a bound along training runs.}

For linear decay the theorem gives about $2DG/\sqrt T$ (\cref{cor:schedules:shapes}), twice the averaging bound. A different argument closes the gap.

\begin{proposition}[Linear decay matches averaging]\label[proposition]{prop:schedules:linear}\index{linear decay}
Let \cref{ass:schedules:convex} hold, let $\gamma>0$, and run \eqref{eq:schedules:sgd} with $\eta_t=\gamma\,(T-t)/T$ for $t=1,\dots,T-1$. Then
\[
  \E\Loss(\theta_T)-\Loss^\star\le\frac{D^2}{2\gamma T}+\frac{\gamma G^2}2 ,
\]
which is $DG/\sqrt T$ for $\gamma=D/(G\sqrt T)$, the bound for the average iterate of stochastic gradient descent with the constant step $\gamma$.
\end{proposition}

\begin{proof}
The plan is to write the last iterate as the average of a sequence that takes constant steps along the same stochastic gradients, and to show by convexity that the regret of that sequence controls the last iterate. Let $z_1=\theta_1$ and $z_{t+1}=z_t-\gamma g_t$ for $t=1,\dots,T$, with $g_t$ the stochastic gradient at $\theta_t$.

\emph{The last iterate is an average.} We claim that $T\theta_t=\sum_{s<t}z_s+(T-t+1)z_t$ for $1\le t\le T$. It holds for $t=1$, and if it holds for $t$, the claimed values at $t+1$ and $t$ differ by $T^{-1}\bigl(z_t+(T-t)z_{t+1}-(T-t+1)z_t\bigr)=\frac{T-t}T(z_{t+1}-z_t)=-\eta_tg_t$, which is the step \eqref{eq:schedules:sgd}. In particular $\theta_T=T^{-1}\sum_{s=1}^Tz_s$.

\emph{A convexity inequality.} Let $p_t=t^{-1}\sum_{s\le t}z_s$, so that $p_1=\theta_1$ and $p_T=\theta_T$. For $2\le t\le T$ the claim reads $T\theta_t=(t-1)p_{t-1}+(T-t+1)z_t$; eliminating $z_t$ with $tp_t=(t-1)p_{t-1}+z_t$ gives
\[
  z_t-\theta_t=\frac{t-1}{T-t+1}\,(\theta_t-p_{t-1}),\qquad p_t=\frac{(t-1)(T-t)\,p_{t-1}+T\theta_t}{t\,(T-t+1)} .
\]
The second identity makes $p_t$ a convex combination of $p_{t-1}$ and $\theta_t$, so with $F_1=\Loss(\theta_1)$ and $F_t=\bigl((t-1)(T-t)F_{t-1}+T\Loss(\theta_t)\bigr)/(t(T-t+1))$, convexity and induction give $\Loss(p_t)\le F_t$; note that $F_T=\Loss(\theta_T)$. Let $s_t\in\partial\Loss(\theta_t)$ and $u\in\R^d$. The first identity and the subgradient inequality at $\theta_t$, applied to $u$ and to $p_{t-1}$, give
\[
  \ip{s_t,z_t-u}\ge\Loss(\theta_t)-\Loss(u)+\frac{t-1}{T-t+1}\bigl(\Loss(\theta_t)-F_{t-1}\bigr)=tF_t-(t-1)F_{t-1}-\Loss(u),
\]
the equality being the definition of $F_t$ rearranged (for $t=1$ read $(t-1)F_{t-1}=0$). The right side telescopes, so $\sum_{t=1}^{T}\ip{s_t,z_t-u}\ge T\bigl(\Loss(\theta_T)-\Loss(u)\bigr)$.

\emph{Conclusion.} Take $u=\theta^\star$ and $s_t=\E[g_t\mid\cF_t]$. Since $z_t$ is $\cF_t$-measurable, $\E\ip{g_t,z_t-\theta^\star}=\E\ip{s_t,z_t-\theta^\star}$. Since $z_{t+1}=z_t-\gamma g_t$, the expansion in the proof of \cref{lem:schedules:anchor}, with $z_t$ in place of $\theta_t$ (it uses only the form of the step, not where $g_t$ was computed), gives $\sum_{t=1}^T\E\ip{g_t,z_t-\theta^\star}\le D^2/(2\gamma)+\gamma TG^2/2$. Divide by $T$.
\end{proof}

The gradient $g_t$ enters $\theta_T$ with the weight $\gamma(T-t)/T$, exactly its total weight in the average of $z_1,\dots,z_T$: the last iterate of linear decay \emph{is} an average, computed without storing one. \Textcite{DefazioEtAl2023} derive the schedule from a conversion of this kind with general weights (Exercise~\ref{ex:schedules:weights}).

\begin{corollary}[Constant steps and linear decay]\label[corollary]{cor:schedules:shapes}
Under the hypotheses of \cref{thm:schedules:lastiterate}, let $H_n=\sum_{j=1}^n1/j$.
\begin{enumerate}
\item For $\eta_t=\eta$ the right side of \eqref{eq:schedules:lastiterate} is $D^2/(2\eta T)+\frac12G^2\eta\,(1+H_{T-1})$. Its minimum over $\eta$ is $DG\sqrt{(1+H_{T-1})/T}$; the step $\eta=D/(G\sqrt T)$ gives $(DG/\sqrt T)(1+\frac12H_{T-1})$.
\item For $\eta_t=\eta\,(T+1-t)/(T+1)$ it is $D^2/(\eta T)+G^2\eta\,(T-\frac12+H_{T-1})/(T+1)$, with minimum $(2DG/\sqrt T)(1+O(T^{-1}\log T))$.
\end{enumerate}
\end{corollary}

\begin{proof}
(i) Here $V_k/R_k=\eta$ and $\eta_k/R_{k+1}=1/(T-k)$, so the sum is $\eta H_{T-1}$, and the first term is $D^2/(2\eta T)+G^2\eta/2$; the minimum of $a/\eta+b\eta$ is $2\sqrt{ab}$. (ii) Put $m=T+1-k$ and $\delta=\eta/(T+1)$. Then $\eta_k=\delta m$, $R_k=\delta m(m+1)/2$, $V_k=\delta^2m(m+1)(2m+1)/6$ and $R_{k+1}=\delta(m-1)m/2$, so the $k$th term of the sum is $\frac{2\delta}{3}\cdot\frac{2m+1}{m-1}$, and $\sum_{m=2}^{T}\frac{2m+1}{m-1}=2(T-1)+3H_{T-1}$. With $R_1=\eta T/2$ and $V_1=\eta^2T(2T+1)/(6(T+1))$, collecting terms gives the formula.
\end{proof}

With the best constant step the bound exceeds the averaging bound by the factor $\sqrt{1+H_{T-1}}\approx\sqrt{\ln T}$, and with the step $D/(G\sqrt T)$ by about $\frac12\ln T$. Its constant is not sharp: \textcite{ZamaniGlineur2025} computed the exact worst case of the last iterate of the subgradient method with constant steps, and for the best constant step it is at most $DG\sqrt{1+\log(N)/4}/\sqrt{N+1}$ after $N$ steps.\sidenote{\Textcite{ZamaniGlineur2025} also show that linear decay attains the lower bound $DG/\sqrt{N+1}$ that holds for all first-order methods, and that no single sequence of step sizes attains it at every $N$: the optimal schedule must know its horizon, which is \cref{obs:schedules:horizon} in the worst case. For $\eta_t\propto1/\sqrt t$, which does not, the last iterate loses a logarithmic factor that cannot be removed \parencite{HarveyEtAl2019}.} For other shapes the sums in \eqref{eq:schedules:lastiterate} are Riemann sums. For a shape $\varphi$ of bounded variation on $[0,1]$ and positive on $[0,1)$, put $\Phi(s)=\int_s^1\varphi$, $\Psi(s)=\int_s^1\varphi^2$ and
\begin{equation}\label{eq:schedules:J}
  J[\varphi]=\frac{\Psi(0)}{\Phi(0)}+\int_0^1\frac{\varphi(s)\,\Psi(s)}{\Phi(s)^2}\,ds .
\end{equation}

\begin{proposition}[The bound for a shape]\label[proposition]{prop:schedules:shape}
Let $\varphi$ be a shape as above such that $\varphi(s)/(1-s)^p$ has a positive limit as $s\to1$, for some $p>0$. Let $\eta_t=\eta\,\varphi(t/(T+1))$ for $t\le T$, with $\eta=\gamma D/(G\sqrt T)$ for a fixed $\gamma>0$. Then $J[\varphi]<\infty$; as $T\to\infty$ the right side of \eqref{eq:schedules:lastiterate} is $(DG/\sqrt T)\bigl(1/(2\gamma\Phi(0))+\gamma J[\varphi]/2+o(1)\bigr)$, and its minimum over $\gamma$ is $DG\sqrt{J[\varphi]/(\Phi(0)T)}\,(1+o(1))$.
\end{proposition}

\begin{proofsketch}
Near $s=1$ the hypothesis gives $\Phi\asymp(1-s)^{p+1}$, $\Psi\asymp(1-s)^{2p+1}$ and $\varphi\Psi/\Phi^2\asymp(1-s)^{p-1}$, which is integrable, so $J[\varphi]<\infty$. With $s=k/(T+1)$, the sums $R_k$ and $V_k$ are Riemann sums of $\eta T\Phi(s)$ and $\eta^2T\Psi(s)$, so for $1-s\gg1/T$ the $k$th term of the sum in \eqref{eq:schedules:lastiterate} is $(\eta/T)\,\varphi\Psi/\Phi^2$ at $s$ up to a factor $1+o(1)$: the sum is a Riemann sum of the integral in \eqref{eq:schedules:J}. The power law at $s=1$ makes the terms with $1-s\le\eps$ contribute $O(\eta\eps^p)$ in total, uniformly in $T$; we omit this estimate, which for the shapes of \cref{tab:schedules:bound} can be done by exact summation as in \cref{cor:schedules:shapes}(ii). The hypothesis cannot be dropped: if $\varphi$ falls faster than every power of $1-s$, as $e^{-0.02/(1-s)}$ does, the last term alone is close to $\eta$, and the bound converges to the formula with $J[\varphi]+1$ in place of $J[\varphi]$. At $T=10^5$ the discrete bound agrees with the limit within $0.02\%$ for the shapes of \cref{tab:schedules:bound}.
\end{proofsketch}

\begin{table}[t]
\begin{fullwidth}
\centering
\small
\begin{tabular}{@{}lccl@{}}
\toprule
shape & $\Phi(0)$ & $J[\varphi]$ & best bound $\times\sqrt T/(DG)$\\
\midrule
constant & $1$ & $1+H_{T-1}$ (exact) & $\sqrt{1+H_{T-1}}$, which is $3.28$ at $T=10^4$\\
linear decay to zero & $1/2$ & $2$ & $2$\\
cosine to zero & $1/2$ & $2.122$ & $2.060$\\
cosine to $\rho=0.1$ & $0.55$ & $1.775+0.1\ln T$ & $2.21$ at $T=10^4$; $2.31$ at $T=10^5$\\
WSD, linear cooldown of fraction $c$ & $1-c/2$ & $2+\ln\frac{2-c}c$ & $2.28$, $2.16$, $2.03$ for $c=0.1$, $0.2$, $0.5$\\
WSD, $1-\sqrt\cdot$ cooldown, $c=0.2$ & $0.867$ & $3.92$ & $2.13$\\
\bottomrule
\end{tabular}
\caption{The last-iterate bound \eqref{eq:schedules:lastiterate} at its best peak (\cref{prop:schedules:shape}); the constant and floor rows are exact discrete values, whose $J$ grows with $T$. A floor $\rho$ adds $\rho\ln T$ to $J$ (Exercise~\ref{ex:schedules:floor}).}\label{tab:schedules:bound}
\end{fullwidth}
\end{table}

\begin{corollary}[The cooldown in the bound]\label[corollary]{cor:schedules:cooldown}
For WSD with peak $\eta$ and a linear cooldown to zero over the last fraction $c\in(0,1]$ of the steps, the right side of \eqref{eq:schedules:lastiterate} is, as $T\to\infty$ with $\eta\sqrt T$ fixed, $U(c)\,(1+o(1))$ with
\begin{equation}\label{eq:schedules:wsdbound}
  U(c)=\frac{D^2}{(2-c)\,\eta T}+\frac{G^2\eta}2\Bigl(2+\ln\frac{2-c}{c}\Bigr).
\end{equation}
\begin{enumerate}
\item With the peak tuned for each $c$, the bound is
\[
  \min_{\eta>0}U(c)=DG\,\sqrt{\frac{2\,\bigl(2+\ln\frac{2-c}{c}\bigr)}{(2-c)\,T}} ,
\]
which decreases in $c$: the bound prefers to decay over the whole run.
\item With the peak fixed, let $\mu=G\eta\sqrt T/D$ be the peak in units of the best peak for linear decay. Then $U$ is minimized at $c^\star=\min\{1,\,2\mu^2/(1+\mu^2)\}$: a peak below the best one asks for a shorter cooldown.
\end{enumerate}
\end{corollary}

\begin{proof}
The plan is to compute $J$ for this shape and apply \cref{prop:schedules:shape}. Here $\Phi(0)=1-c/2$ and $\Psi(0)=1-2c/3$. On the cooldown $\Phi(s)=(1-s)^2/(2c)$ and $\Psi(s)=(1-s)^3/(3c^2)$, so the integrand of \eqref{eq:schedules:J} is $4/(3c)$ and contributes $4/3$. On the stable phase, with $w=1-s-c/2$, the integrand is $(w-c/6)/w^2$, with integral $\ln\frac{2-c}c-\frac13+\frac{c}{3(2-c)}$ over $w\in[c/2,1-c/2]$. Adding $\Psi(0)/\Phi(0)=\frac{2(3-2c)}{3(2-c)}$ gives $J=2+\ln\frac{2-c}c$.
(i) The tuned value is $DG\sqrt{2F(c)/T}$ with $F(c)=(2+\ln r)/(2-c)$ and $r=(2-c)/c\ge1$, and $F'(c)=(1-r+\ln r)/(2-c)^2\le0$ because $\ln r\le r-1$.
(ii) The derivative $U'(c)=D^2/((2-c)^2\eta T)-G^2\eta/(c(2-c))$, multiplied by $c(2-c)^2\eta T/D^2>0$, is $c(1+\mu^2)-2\mu^2$, which changes sign from negative to positive at $2\mu^2/(1+\mu^2)$; if that exceeds $1$, then $U$ decreases on $(0,1]$.
\end{proof}

So the bound reproduces \cref{obs:schedules:cooldown} in a weak form: a cooldown removes the logarithm of the constant schedule, and its benefit saturates. It does not single out $20\%$; with tuned peaks it prefers a decay over the whole run, as in the experiments of \textcite{DefazioEtAl2023} and \textcite{BergsmaEtAl2025}, and it prefers the $1-\sqrt\cdot$ cooldown of \textcite{HageleEtAl2024} to a linear one.

\section{The theory: the river valley}\label{sec:schedules:river}

In the river valley the expected loss can be computed exactly for every schedule. For a hill $i$ write $q_{i,t}=(1-\eta_th_i)^2$ and $P_i(k,T)=\prod_{t=k}^{T-1}q_{i,t}$ (an empty product is $1$), the contraction of its second moment between steps $k$ and $T$. When $0\le\eta h_i<2$, let
\[
  v_i(\eta)=\frac{\eta\,\sigma_i^2}{h_i\,(2-\eta h_i)} ,
\]
the stationary second moment of the hill at the constant step $\eta$ (\cref{lem:toolkit:lyapunov}).

\begin{theorem}[Stochastic gradient descent in the river valley]\label{thm:schedules:river}
Let \cref{ass:schedules:river} hold and let $(\eta_t)$ be a schedule. For each hill $i$ and each $T\ge1$:
\begin{enumerate}
\item $\E y_{T,i}^2=P_i(1,T)\,y_{1,i}^2+\sigma_i^2\sum_{k=1}^{T-1}\eta_k^2\,P_i(k+1,T)$.
\item If $0\le\eta_th_i<2$ for $t<T$, then
\begin{multline}\label{eq:schedules:annealing}
  \E y_{T,i}^2=v_i(\eta_1)+P_i(1,T)\bigl(y_{1,i}^2-v_i(\eta_1)\bigr)\\
  -\sum_{k=2}^{T-1}\bigl(v_i(\eta_{k-1})-v_i(\eta_k)\bigr)\bigl(1-P_i(k,T)\bigr).
\end{multline}
\item If $0\le\eta_th_i\le\frac12$ for $k\le t<T$, then $e^{-2h_i(A_T-A_k)-2h_i^2\sum_{t=k}^{T-1}\eta_t^2}\le P_i(k,T)\le e^{-2h_i(A_T-A_k)}$.
\item If the river is linear, then $\E r(x_T)=r(x_1)-\norm{\nabla r}_2^2\,A_T$. If it is quadratic, each of its eigendirections satisfies (i)--(iii) with its own curvature and noise.
\end{enumerate}
Consequently $\E\Loss(\theta_T)=\E r(x_T)+\frac12\sum_ih_i\,\E y_{T,i}^2$ is given in closed form by the schedule.
\end{theorem}

\begin{proof}
The plan is to derive the second-moment recursion of each coordinate and to rewrite it around the moving equilibrium $v_i(\eta_t)$. (i) The hill coordinate moves by $y_{t+1,i}=(1-\eta_th_i)y_{t,i}-\eta_t\zeta_{t,i}$. Square it and condition on $\cF_t$; the cross term vanishes because $y_{t,i}$ is $\cF_t$-measurable and $\E[\zeta_{t,i}\mid\cF_t]=0$, so $\E[y_{t+1,i}^2\mid\cF_t]=q_{i,t}y_{t,i}^2+\eta_t^2\sigma_i^2$. Take expectations and unroll from $t=1$.
(ii) Write $u_t=\E y_{t,i}^2$ and $v_t=v_i(\eta_t)$; the step size $\eta_T$ does not affect $y_{T,i}$, so we may take $\eta_T=\eta_{T-1}$ and $v_T$ is defined. A direct computation gives $q_{i,t}v_t+\eta_t^2\sigma_i^2=v_t$, so subtracting it from the recursion of (i) gives $u_{t+1}-v_t=q_{i,t}(u_t-v_t)$, that is, $u_{t+1}-v_{t+1}=q_{i,t}(u_t-v_t)+(v_t-v_{t+1})$. Unrolling, $u_T-v_T=P_i(1,T)(u_1-v_1)+\sum_{k=2}^{T}(v_{k-1}-v_k)P_i(k,T)$. The term $k=T$ has $P_i(T,T)=1$ and cancels the unused $v_T$, and $v_{T-1}=v_1-\sum_{k=2}^{T-1}(v_{k-1}-v_k)$ gives \eqref{eq:schedules:annealing}.
(iii) For $0\le z\le\frac12$, $e^{-z-z^2}\le1-z\le e^{-z}$: the right inequality is convexity of the exponential, and $\ln(1-z)+z+z^2$ vanishes at $0$ and has derivative $z(1-2z)/(1-z)\ge0$. Apply it with $z=\eta_th_i$, square, and multiply.
(iv) A linear river has $x_{t+1}=x_t-\eta_t(\nabla r+\zeta_{t,x})$, so $\E x_T=x_1-A_T\nabla r$, and $r$ is affine. The eigendirections of a quadratic river are decoupled coordinates of the same form as the hills.
\end{proof}

Part (ii), the \term{annealing decomposition}\index{annealing decomposition}, is the structure behind the chapter: the second moment of a hill at step $T$ is its equilibrium at the initial step size, plus the fading memory of the start, minus one term for each decrease of the step size, the decrease of the equilibrium times the fraction $1-P_i(k,T)$ of it realized since. By (iii) that fraction depends on the area since the decrease, $1-e^{-2h_i(A_T-A_k)}$ up to relative corrections of order $\eta h_i$; an increase of the step size, as in warmup, enters with the opposite sign.

\begin{corollary}[Noise floor and relaxation]\label[corollary]{cor:schedules:floor}
Under \cref{ass:schedules:river}, let $0<\eta h_i<2$ for every hill.
\begin{enumerate}
\item At the constant step $\eta$ the hill part of the expected loss converges, at the rate $(1-\eta h_i)^2$ per step in hill $i$, to the \term{noise floor}\index{noise floor} $\Loss_{\mathrm{hill}}(\eta)=\sum_{i}\eta\sigma_i^2/(2(2-\eta h_i))$, which lies between $\frac\eta4\sum_i\sigma_i^2$ and $\eta\sum_i\sigma_i^2/(4-2\eta\max_ih_i)$.
\item If $\E y_{k,i}^2=v_i(\eta)$ and the step size is $\eta'\in[0,\eta)$ from step $k$ on, then $\E y_{T,i}^2=v_i(\eta')+\bigl(v_i(\eta)-v_i(\eta')\bigr)(1-\eta'h_i)^{2(T-k)}$ for $T\ge k$.
\end{enumerate}
\end{corollary}

\begin{proof}
(i) In \eqref{eq:schedules:annealing} all differences vanish and $P_i(1,T)=(1-\eta h_i)^{2(T-1)}$; multiply by $h_i/2$ and sum, and use $2-\eta\max_ih_i\le2-\eta h_i\le2$. (ii) The recursion $u_{t+1}-v_i(\eta')=(1-\eta'h_i)^2(u_t-v_i(\eta'))$ of the proof of \cref{thm:schedules:river}(ii) holds for $t\ge k$.
\end{proof}

So at a constant learning rate each hill adds about $\eta\sigma_i^2/4$ to the loss, a floor proportional to the learning rate and to the noise, hence to $1/B$. A drop of the learning rate removes the excess on the \term{relaxation time}\index{relaxation time} $1/(\eta'h_i(2-\eta'h_i))\approx1/(2\eta'h_i)$ of the \emph{new} step size, so a decay to zero never quite finishes.

\begin{heuristic}[The cooldown law]\index{cooldown!law}
Take one hill in equilibrium at the peak $\eta$, and a linear cooldown to zero over the last $K$ steps. Measure time by the area: to first order in $\eta h$ the recursion of \cref{thm:schedules:river}(i) is $dv/dA=-2hv+\eta(A)\sigma^2$, because a step changes $v$ by $-(2\eta_th-\eta_t^2h^2)v+\eta_t^2\sigma^2$ and the area by $\eta_t$. The cooldown has area $\eta K/2$, and with $\varrho$ the area that remains before the end the step size is $\eta(\varrho)=\sqrt{2\eta\varrho/K}$. Solving the linear equation from $v=\eta\sigma^2/(2h)$,
\[
  v_{\mathrm{end}}=\frac{\eta\sigma^2}{2h}e^{-\eta hK}+\sigma^2\int_0^{\eta K/2}e^{-2h\varrho}\sqrt{2\eta\varrho/K}\,d\varrho ,
\]
an incomplete gamma function. With $\kappa=\eta hK$ the two terms combine into
\begin{equation}\label{eq:schedules:cooldownlaw}
  \frac h2\,v_{\mathrm{end}}=\frac{\eta\sigma^2}4\,\chi(\kappa),\qquad \chi(\kappa)=\frac{\sqrt\pi}2\,\frac{\erf\sqrt\kappa}{\sqrt\kappa} .
\end{equation}
A short cooldown has $\chi(\kappa)=1-\kappa/3+O(\kappa^2)$ and removes almost nothing; a long one leaves $\chi(\kappa)\approx\frac{\sqrt\pi}2\kappa^{-1/2}$. The residual noise after a linear cooldown of $K$ steps decays like $K^{-1/2}$, not exponentially: the noise injected within the last area $1/(2h)$, at step sizes of order $(\eta/(hK))^{1/2}$, has no area left in which to be forgotten.\qedhere
\end{heuristic}

The exact recursion confirms \eqref{eq:schedules:cooldownlaw} within relative errors of order $\eta h$: started at the equilibrium $v(\eta)$, with the steps $\eta(1-j/K)$ for $j=1,\dots,K$, for $\eta h=0.01$ and $K=100$, $1000$, $10^4$ it leaves $0.746$, $0.279$ and $0.0882$ of the floor, against $\chi=0.747$, $0.280$ and $0.0886$. The cooldown that removes the most noise in $K$ steps, ignoring the river, is a $1/t$ decay, with a residual of order $1/K$ (Exercise~\ref{ex:schedules:greedy}).\sidenote{The error function enters as $\int_0^x\sqrt u\,e^{-u}du=\frac{\sqrt\pi}2\erf\sqrt x-\sqrt x\,e^{-x}$. A cooldown whose step size vanishes like the $p$th power of the remaining steps leaves a tail $\kappa^{-p/(p+1)}$: $\kappa^{-1/2}$ for the linear and $1-\sqrt\cdot$ shapes, $\kappa^{-2/3}$ for a cosine.}

\begin{heuristic}[The best cooldown]\index{cooldown!optimal length}
Let the river be linear, so that by \cref{thm:schedules:river}(iv) it loses $\omega=\norm{\nabla r}_2^2$ per unit of area, and replace the last $K$ steps of a constant schedule by a linear cooldown. The cooldown gives up the area $\eta K/2$ and removes noise by \eqref{eq:schedules:cooldownlaw}, so the final loss is, up to a constant, $\omega\eta K/2+\sum_i\frac14\eta\sigma_i^2\chi(\eta h_iK)$. For one hill with $\eta hK\gg1$ the derivative vanishes at
\begin{equation}\label{eq:schedules:Kstar}
  K^\star=\Bigl(\frac{\sqrt\pi\,\sigma^2}{8\,\omega}\Bigr)^{2/3}(\eta h)^{-1/3},
\end{equation}
where the residual noise is twice the river progress given up. The best cooldown is a fixed number of steps, and its \emph{fraction} of the run falls like $1/T$. A river that flattens changes this. If the river loss is a power of the area, $\bar r(A)=L_\infty+CA^{-\alpha}$, then its slope at the end of a run of area $\eta T$ is $\alpha C(\eta T)^{-1-\alpha}$, and the same balance, with $1-c/2\approx1$, gives the best cooldown fraction
\begin{equation}\label{eq:schedules:cstar}
  c^\star\approx\Bigl(\frac{\sqrt\pi}{8\alpha C}\sum_i\frac{\sigma_i^2}{\sqrt{h_i}}\Bigr)^{2/3}\,\eta^{(1+2\alpha)/3}\,T^{(2\alpha-1)/3} .
\end{equation}
A fixed fraction is optimal at every horizon exactly when $\alpha=\frac12$; a river that falls faster than the inverse square root of the area asks for a longer cooldown in a longer run, a slower one for a shorter cooldown.\qedhere
\end{heuristic}

For one hill with $\eta h=0.01$, $\sigma^2=1$ and $\omega=10^{-4}$, the exact recursion puts the best linear cooldown at $790$ steps, \eqref{eq:schedules:Kstar} at $789$; the asymptotic form fails, as it must, when $\eta hK^\star$ is of order one. The decomposition also explains the laws of \cref{obs:schedules:laws}.

\begin{heuristic}[From the river valley to the empirical laws]
Multiply \eqref{eq:schedules:annealing} by $h_i/2$, sum over the hills, and keep the first order in $\eta h_i$, where $\frac{h_i}2v_i(\eta)\approx\frac14\eta\sigma_i^2$ and $1-P_i(k,T)\approx1-e^{-2h_i(A_T-A_k)}$. Let $\schedulesG(u)=\sum_i\frac14\sigma_i^2\bigl(1-e^{-2h_iu}\bigr)$, a saturation function built from the Laplace transform of the noise-weighted hill spectrum. Once the initial transient has faded,
\begin{equation}\label{eq:schedules:law}
  \E\Loss(\theta_T)\approx\E r(x_T)+\frac{\eta_1}4\sum_i\sigma_i^2-\sum_{k=2}^{T-1}(\eta_{k-1}-\eta_k)\,\schedulesG(A_T-A_k).
\end{equation}
This is the form of \cref{obs:schedules:laws}: an area term, close to $L_\infty+CA_T^{-\alpha}$ for a quadratic river with a power-law spectrum, and a loss drop that adds, for each decrease of the learning rate, the decrease times a saturating function of the area since. If the noise is spread over many hills and its distribution over their curvatures, with weights $\sigma_i^2$, is a gamma distribution with shape $\nu$ and scale $\bar h$, then $\schedulesG(u)=\frac14\bigl(\sum_i\sigma_i^2\bigr)\bigl(1-(1+2\bar hu)^{-\nu}\bigr)$, the power-law saturation of \textcite{LuoEtAl2025}. A single hill at a nearly constant step $\bar\eta$ after the decreases gives the saturation $1-e^{-2\bar\eta h(T-k)}$, counted in steps, which is the annealing area of \textcite{TissueEtAl2025}. Finally, the continuum form of \cref{thm:schedules:river}(i) writes the noise term as $\sum_k\eta_k^2\schedulesK(A_T-A_{k+1})$ with the kernel $\schedulesK=\schedulesG'$, a convolution of the squared schedule with a kernel in area-time; \textcite{LiEtAl2025} derive a functional scaling law of this kind for stochastic gradient descent on kernel regression with power-law spectra.\qedhere
\end{heuristic}

The model predicts what the laws leave open: the loss relaxes after a decrease in \emph{area}, not in steps, because a quadratic hill contracts by $(1-\eta h)^2\approx e^{-2h\eta}$ per step, that is, by $e^{-2h}$ per unit of area. \Textcite{TissueEtAl2025} count the relaxation in steps, \textcite{LuoEtAl2025} in area with a fitted correction that depends on the learning rate, which measures a departure from a fixed quadratic landscape.

\begin{remark}[Branching, rewarming and averaging]\label[remark]{rem:schedules:wsd}
After a long stable phase the hills have forgotten everything but the peak, so by \eqref{eq:schedules:annealing} the loss after a cooldown depends on the stable phase only through its area (\cref{obs:schedules:horizon}); rewarming a cooled checkpoint refills the hills within a relaxation time and keeps the river progress. Averaging can replace the decay, as in schedule-free methods \parencite{DefazioEtAl2024} (\cref{ch:momentum}): the mean of the last $K$ iterates at a constant step gives up the same area as a linear cooldown of $K$ steps and leaves a residual of order $1/(\eta hK)$, not $(\eta hK)^{-1/2}$ (Exercise~\ref{ex:schedules:averaging}).
\end{remark}

\section{Reading the theory}\label{sec:schedules:reading}

\paragraph{Loss curves.} \Cref{fig:schedules:river} runs the four schedules in a river valley built to look like a language model: a quadratic river in $200$ directions whose loss falls like the inverse square root of the area, and $40$ noisy hills with curvatures between $0.2$ and $2$. The points, means of $100$ runs, agree with the exact curves within their sampling error (final values $0.203$, $0.156$, $0.158$, $0.131$ against $0.207$, $0.157$, $0.158$, $0.131$). The constant schedule reaches the floor $\Loss_{\mathrm{hill}}(0.04)=0.102$ within a few hundred steps and follows the river down at that height. In the cooldown of WSD the hill part falls from $0.102$ to $0.021$ while the river part continues from $0.117$ to $0.110$: four fifths of the floor removed for a tenth of the area. Cosine and linear decay remove almost all of the floor but cover half the area, and their river part ends at $0.149$. At a common peak the model ranks WSD first, cosine and linear decay together, the constant schedule last.

\begin{figure}[t]
\begin{fullwidth}
\centering
\begin{tikzpicture}
\begin{groupplot}[group style={group size=2 by 1,horizontal sep=17mm},wplot,width=82mm,height=48mm,xmin=0,xmax=4000,xtick={0,1000,2000,3000,4000},scaled x ticks=false,xticklabel style={/pgf/number format/1000 sep={}}]
\nextgroupplot[xlabel={step $t$},ylabel={expected loss},ymode=log,ymin=0.1,ymax=0.65,ytick={0.1,0.15,0.2,0.3,0.4,0.6},yticklabels={$0.1$,$0.15$,$0.2$,$0.3$,$0.4$,$0.6$},legend style={at={(0.5,1.03)},anchor=south},legend columns=4]
\addplot[theory] table[x=t,y=const] {figures/data/12-schedules-river.dat};
\addplot[theory3] table[x=t,y=cosine] {figures/data/12-schedules-river.dat};
\addplot[theory2] table[x=t,y=linear] {figures/data/12-schedules-river.dat};
\addplot[muted,line width=1.0pt,dashdotted] table[x=t,y=wsd] {figures/data/12-schedules-river.dat};
\addplot[only marks,mark=*,mark size=1.1pt,sea,each nth point=2,filter discard warning=false] table[x=t,y=const] {figures/data/12-schedules-river-sgd.dat};
\addplot[only marks,mark=*,mark size=1.1pt,sandline!55!black,each nth point=2,filter discard warning=false] table[x=t,y=cosine] {figures/data/12-schedules-river-sgd.dat};
\addplot[only marks,mark=square*,mark size=1.0pt,seadeep,each nth point=2,filter discard warning=false] table[x=t,y=linear] {figures/data/12-schedules-river-sgd.dat};
\addplot[only marks,mark=triangle*,mark size=1.3pt,muted,each nth point=2,filter discard warning=false] table[x=t,y=wsd] {figures/data/12-schedules-river-sgd.dat};
\legend{constant,cosine,linear,WSD}
\nextgroupplot[xlabel={step $t$},ylabel={river and hill parts},ymode=log,ymin=0.004,ymax=0.65,legend style={at={(0.5,1.03)},anchor=south},legend columns=2]
\addplot[muted,line width=1.0pt,dashdotted] table[x=t,y=wsdriver] {figures/data/12-schedules-river.dat};
\addplot[muted,line width=0.8pt,densely dotted] table[x=t,y=wsdhill] {figures/data/12-schedules-river.dat};
\addplot[theory3] table[x=t,y=cosriver] {figures/data/12-schedules-river.dat};
\addplot[sandline!55!black,line width=0.8pt,densely dotted] table[x=t,y=coshill] {figures/data/12-schedules-river.dat};
\addplot[only marks,mark=triangle*,mark size=1.3pt,muted,each nth point=2,filter discard warning=false] table[x=t,y=wsdriver] {figures/data/12-schedules-river-sgd.dat};
\addplot[only marks,mark=triangle*,mark size=1.3pt,muted,each nth point=2,filter discard warning=false] table[x=t,y=wsdhill] {figures/data/12-schedules-river-sgd.dat};
\addplot[only marks,mark=*,mark size=1.1pt,sandline!55!black,each nth point=2,filter discard warning=false] table[x=t,y=cosriver] {figures/data/12-schedules-river-sgd.dat};
\addplot[only marks,mark=*,mark size=1.1pt,sandline!55!black,each nth point=2,filter discard warning=false] table[x=t,y=coshill] {figures/data/12-schedules-river-sgd.dat};
\legend{WSD river,WSD hills,cosine river,cosine hills}
\end{groupplot}
\end{tikzpicture}
\end{fullwidth}
\caption{Four schedules in a river valley: curves are the exact expected loss (\cref{thm:schedules:river}), points the mean of $100$ runs of stochastic gradient descent with Gaussian noise (seed $12$), every $200$ steps. River: $200$ directions with curvatures $0.02\,j^{-2}$, initial losses proportional to $j^{-2}$ (total $0.5$), no noise. Hills: $40$, curvatures log-spaced in $[0.2,2]$, started at $0$, noise variance $1/4$ each. Peak $0.04$, $T=4000$, no warmup; WSD cools down linearly over the last $800$ steps. Left: total loss. Right: river part (upper curves) and hill part (lower curves). The points agree with the curves within sampling error (at most $2.7\%$ apart, less than three standard errors); the river carries no noise, so its points coincide with its curves. Script: \texttt{12-schedules-river.py}.}
\label{fig:schedules:river}
\end{figure}

That ranking is not a prediction about tuned runs. With each peak tuned separately among the peaks at which every step is stable, $\eta\max_ih_i<2$, the model raises the peaks of the decaying schedules far above the common one (cosine to that limit, linear decay to $0.64$, WSD to $0.21$) and ranks cosine first ($0.052$, against $0.074$ for linear decay, $0.096$ for WSD and $0.195$ for constant steps): a larger peak buys area, and a long decay removes the larger floor. Cosine even gains from peaks up to nearly twice that limit ($0.050$ at the peak $1.9$): the stiff hills grow during its first, unstable steps, and the contraction during the later steps undoes the growth, which only exactly quadratic hills allow. Nothing in the model penalizes a large peak short of divergence (\cref{ch:eos,ch:spikes}). What is robust are the decomposition and its time scales: the cooldown drop is the noise floor at the peak, about $\eta\tr(\Sigma_{\mathrm{hill}})/(4B)$, so it should double when the peak doubles or the batch size halves; it takes a few relaxation times of the slowest noisy direction; and it leaves a residual decaying only like $(\eta hK)^{-1/2}$. The curves are exact for this model, so their agreement with the points is guaranteed up to sampling error; it must fail once the curvature changes along the run, the noise depends on the position, or the optimizer is not plain stochastic gradient descent.

\paragraph{The convex bound along a run.} \Cref{fig:schedules:bound} (left) evaluates \eqref{eq:schedules:lastiterate} at every step on the \term{cone}\index{cone (test problem)} $\Loss(\theta)=\E_a\abs{\ip{a,\theta-\theta^\star}}=\sqrt{2/\pi}\,\norm{\theta-\theta^\star}_2$ with $a\sim\cN(0,I_{100})$, which is convex with a nonsmooth minimum and has $G^2=d$ exactly. The shapes agree with the simulation and with \cref{obs:schedules:convex}: constant steps plateau, WSD plateaus and then drops, cosine and linear decay descend smoothly, and both rank the final values linear decay, cosine, WSD, constant. The values do not agree: the bound is $7.6$ to $9.5$ times the simulated suboptimality at the end. On the cone the last iterate of constant steps settles exactly at the averaging floor $\eta G^2/2$, with no logarithm, and it forgets a change of the step size within about $130$ steps whatever the step size, so its final loss is set by the last steps. Tuned comparisons (right panel) suffer: the bound charges a $20\%$ cooldown $8\%$ over linear decay, the cone $68\%$; cosine $3\%$, the cone $31\%$. The bound covers every convex Lipschitz function, and the functions on which the last iterate comes within a constant factor of it are specially constructed \parencite{HarveyEtAl2019,ZamaniGlineur2025}; it agrees with training curves in shapes and rankings, not in values.

\begin{figure}[t]
\begin{fullwidth}
\centering
\begin{tikzpicture}
\begin{groupplot}[group style={group size=2 by 1,horizontal sep=17mm},wplot,width=82mm,height=48mm]
\nextgroupplot[xlabel={step $t$},ylabel={$\E\Loss(\theta_t)-\Loss^\star$},ymode=log,xmin=0,xmax=1000,ymin=0.04,ymax=4,legend style={at={(0.5,1.03)},anchor=south},legend columns=4,xticklabel style={/pgf/number format/1000 sep={}}]
\addplot[theory,restrict x to domain=10:1000] table[x=t,y=const] {figures/data/12-schedules-bound.dat};
\addplot[theory3,restrict x to domain=10:1000] table[x=t,y=cosine] {figures/data/12-schedules-bound.dat};
\addplot[theory2,restrict x to domain=10:1000] table[x=t,y=linear] {figures/data/12-schedules-bound.dat};
\addplot[muted,line width=1.0pt,dashdotted,restrict x to domain=10:1000] table[x=t,y=wsd] {figures/data/12-schedules-bound.dat};
\addplot[only marks,mark=*,mark size=1.0pt,sea,restrict x to domain=10:1000,each nth point=2,filter discard warning=false] table[x=t,y=const] {figures/data/12-schedules-bound-sgd.dat};
\addplot[only marks,mark=*,mark size=1.0pt,sandline!55!black,restrict x to domain=10:1000,each nth point=2,filter discard warning=false] table[x=t,y=cosine] {figures/data/12-schedules-bound-sgd.dat};
\addplot[only marks,mark=square*,mark size=0.9pt,seadeep,restrict x to domain=10:1000,each nth point=2,filter discard warning=false] table[x=t,y=linear] {figures/data/12-schedules-bound-sgd.dat};
\addplot[only marks,mark=triangle*,mark size=1.2pt,muted,restrict x to domain=10:1000,each nth point=2,filter discard warning=false] table[x=t,y=wsd] {figures/data/12-schedules-bound-sgd.dat};
\legend{constant,cosine,linear,WSD}
\nextgroupplot[xlabel={cooldown fraction $c$},ylabel={tuned final value / linear decay},xmode=log,xmin=0.02,xmax=1.05,ymin=0.9,ymax=2.2,xtick={0.02,0.05,0.1,0.2,0.5,1},xticklabels={$0.02$,$0.05$,$0.1$,$0.2$,$0.5$,$1$},legend style={at={(0.5,1.03)},anchor=south},legend columns=2]
\addplot[theory] table[x=c,y=boundrel] {figures/data/12-schedules-bound-tunedcurve.dat};
\addplot[simulation] table[x=c,y=sgdrel] {figures/data/12-schedules-bound-tuned.dat};
\addplot[guide,domain=0.02:1.05] {1};
\legend{bound \eqref{eq:schedules:lastiterate},SGD on the cone}
\end{groupplot}
\end{tikzpicture}
\end{fullwidth}
\caption{The last-iterate bound against stochastic gradient descent on the cone $\sqrt{2/\pi}\,\norm{\theta-\theta^\star}_2$ in $d=100$ dimensions, $D=1$, $G^2=d$. Left: the bound \eqref{eq:schedules:lastiterate} for $\theta_t$, computed from $\eta_1,\dots,\eta_t$ (curves), and the mean of $200$ runs (points, seed $21$), $T=1000$, common peak $D/(G\sqrt T)$; WSD cools down over the last $200$ steps. At $t=T$ the bound is $7.6$ to $9.5$ times the simulated value. Right: WSD with cooldown fraction $c$, peak tuned for each $c$, final value relative to linear decay; the bound is minimized exactly, the simulation takes the best of $17$ peaks between a quarter of and four times $D/(G\sqrt T)$, $40$ runs each (seed $22$). Script: \texttt{12-schedules-bound.py}.}
\label{fig:schedules:bound}
\end{figure}

\paragraph{How long to cool down.} \Cref{fig:schedules:cooldown} tests the cooldown heuristics in the valley of \cref{fig:schedules:river} and in one whose river falls like $A^{-1}$. For $\alpha=\frac12$ the best cooldown fraction is $0.24$, $0.28$ and $0.27$ at $T=1000$, $4000$ and $16000$, fixed as \eqref{eq:schedules:cstar} predicts; for $\alpha=1$ it is $0.21$, $0.35$ and $0.48$, a factor $2.3$ over a sixteenfold increase of $T$, against $16^{1/3}=2.5$. The continuum minimizer coincides with the exact one within the grid spacing, and the minimum is flat, as in \cref{obs:schedules:cooldown}: at $\alpha=\frac12$ and $T=4000$ every fraction between $0.18$ and $0.41$ is within $1\%$ of the best. At fractions this large the peak dependence of \eqref{eq:schedules:cstar} is approximate: doubling the peak multiplies $c^\star$ by $1.42$ to $1.55$ for $\alpha=\frac12$ and by $1.56$ to $1.72$ for $\alpha=1$, against $2^{2/3}\approx1.59$ and $2$. The value $c^\star\approx0.27$ depends on the noise, the curvatures and the river; the scaling with $T$ and its dependence on $\alpha$ are the prediction.

\begin{figure}[t]
\begin{fullwidth}
\centering
\begin{tikzpicture}
\begin{groupplot}[group style={group size=2 by 1,horizontal sep=17mm},wplot,width=82mm,height=48mm,xmode=log,xmin=0.01,xmax=1.05,xtick={0.01,0.03,0.1,0.3,1},xticklabels={$0.01$,$0.03$,$0.1$,$0.3$,$1$},ymode=log]
\nextgroupplot[xlabel={cooldown fraction $c$},ylabel={final expected loss},title={river $\propto A^{-1/2}$},ymin=0.055,ymax=0.36,ytick={0.06,0.1,0.15,0.2,0.3},yticklabels={$0.06$,$0.1$,$0.15$,$0.2$,$0.3$}]
\addplot[theory] table[x=c,y=a05T1000] {figures/data/12-schedules-cooldown.dat};
\addplot[theory3] table[x=c,y=a05T4000] {figures/data/12-schedules-cooldown.dat};
\addplot[theory2] table[x=c,y=a05T16000] {figures/data/12-schedules-cooldown.dat};
\addplot[simulation] table[x=c,y=a05T1000] {figures/data/12-schedules-cooldown-sgd.dat};
\addplot[simulation] table[x=c,y=a05T4000] {figures/data/12-schedules-cooldown-sgd.dat};
\addplot[simulation] table[x=c,y=a05T16000] {figures/data/12-schedules-cooldown-sgd.dat};
\addplot[only marks,mark=diamond*,mark size=2.6pt,ink,mark options={fill=white}] table[x=cstara05,y=lossa05] {figures/data/12-schedules-cooldown-opt.dat};
\nextgroupplot[xlabel={cooldown fraction $c$},ylabel={final expected loss},title={river $\propto A^{-1}$},ymin=0.015,ymax=0.3,ytick={0.02,0.05,0.1,0.2},yticklabels={$0.02$,$0.05$,$0.1$,$0.2$},legend style={at={(0.97,0.63)},anchor=east,font=\scriptsize},legend columns=1]
\addplot[theory] table[x=c,y=a1T1000] {figures/data/12-schedules-cooldown.dat};
\addplot[theory3] table[x=c,y=a1T4000] {figures/data/12-schedules-cooldown.dat};
\addplot[theory2] table[x=c,y=a1T16000] {figures/data/12-schedules-cooldown.dat};
\addplot[simulation] table[x=c,y=a1T1000] {figures/data/12-schedules-cooldown-sgd.dat};
\addplot[simulation] table[x=c,y=a1T4000] {figures/data/12-schedules-cooldown-sgd.dat};
\addplot[simulation] table[x=c,y=a1T16000] {figures/data/12-schedules-cooldown-sgd.dat};
\addplot[only marks,mark=diamond*,mark size=2.6pt,ink,mark options={fill=white}] table[x=cstara1,y=lossa1] {figures/data/12-schedules-cooldown-opt.dat};
\legend{$T=1000$,$T=4000$,$T=16000$}
\end{groupplot}
\end{tikzpicture}
\end{fullwidth}
\caption{The final loss of WSD against its cooldown fraction $c$, in the valley of \cref{fig:schedules:river} (left, river loss $\propto A^{-1/2}$) and in the same valley with initial river losses $\propto j^{-3}$ (right, $\propto A^{-1}$), at the peak $0.04$, for $T=1000$, $4000$, $16000$ (top to bottom). Curves: exact expected loss (\cref{thm:schedules:river}); points: mean of $40$ runs of stochastic gradient descent (seed $31$), within sampling error of the curves (at most $4\%$ apart, about two standard errors); diamonds: minimizer of the continuum approximation, with the river loss evaluated at the area of the schedule and the hills by \eqref{eq:schedules:cooldownlaw}, plotted at the exact loss. Script: \texttt{12-schedules-cooldown.py}.}
\label{fig:schedules:cooldown}
\end{figure}

\section{What the theory misses}\label{sec:schedules:misses}

\paragraph{What was idealized away.} The convex bound holds for every convex Lipschitz loss and so cannot see the scales of a particular one (\cref{fig:schedules:bound}). The river valley keeps two time scales but fixes everything else: quadratic hills of constant curvature, decoupled from the river; additive noise with a fixed covariance, whereas in a network the gradient covariance follows the Hessian and changes along the run; and plain stochastic gradient descent, without momentum, Adam's normalization, clipping or weight decay. The river's exponent $\alpha$ and the spread of the noise over the hills are inputs, not outputs, and they set the constants of every prediction. Neither model needs warmup, which \cref{ch:warmup} explains by the curvature that the early steps meet.

\paragraph{Where it breaks.} (1) In the river valley only stability limits the peak of a long decay, so with tuned peaks the model ranks cosine well ahead of WSD, whereas WSD matches cosine in language models (\cref{obs:schedules:cooldown}); what limits the peak in practice, instabilities and loss spikes, is outside the model (\cref{ch:eos,ch:spikes}). (2) The landscape moves with the schedule. In full-batch training the sharpness rises while it is below $2/\eta$ and then stays near $2/\eta$ \parencite{CohenEtAl2021}, so a decaying learning rate lets the sharpest directions sharpen again, and part of the oscillation that a cooldown removes is the oscillation at the edge of stability (\cref{ch:eos}), whose size is set by the dynamics rather than by noise. The two mechanisms differ in a testable way: the noise floor of \cref{cor:schedules:floor} scales like $1/B$, the oscillation at the edge does not. (3) The empirical laws relax in steps, or in area with a fitted correction, not exactly in area as a fixed quadratic hill does (\cref{sec:schedules:river}). (4) With AdamW the norms of normalized weights equilibrate on the time scale $1/(\eta\lambda)$, which a cooldown lengthens (\cref{ch:decay}), so part of what a cooldown does acts through weight norms that the river valley does not have. (5) The convex bound prefers a decay over the whole run, as do the experiments of \textcite{DefazioEtAl2023} and \textcite{BergsmaEtAl2025}, but it underestimates the penalty of short cooldowns on the cone by an order of magnitude, so its agreement with practice is one of direction. And the river valley is itself a conjecture about language models, supported by a toy derivation and by the laws it implies, not derived from a network \parencite{WenEtAl2025}.

\paragraph{Repairs and open problems.} The empirical laws are repairs of the kind the model suggests: \textcite{TissueEtAl2025} and \textcite{LuoEtAl2025} keep the area term and the annealing structure of \eqref{eq:schedules:law} and fit the saturation, and \textcite{LiEtAl2025} derive a functional law from kernel regression with power-law spectra, a noisy quadratic model whose river and hills are the two ends of one spectrum. \Textcite{DefazioEtAl2023} refine the schedule with observed gradient norms and obtain a warmup and a fast final decay; schedule-free methods \parencite{DefazioEtAl2024} average instead of decaying, and \textcite{HageleEtAl2024} compare weight averaging with cooldowns. For least squares the last iterate under step decay \parencite{GeEtAl2019} and under decaying steps in the overparameterized regime \parencite{WuEtAl2022} is understood. Open: a river valley derived from a network (Exercise~\ref{ex:schedules:derive}); a solved model whose hills sit at the edge of stability; a mechanism for the learning-rate dependence of the relaxation (Exercise~\ref{ex:schedules:relax}); and a derivation of the cooldown fraction near $20\%$ that practitioners use.

\section{Back in the wild}\label{sec:schedules:back}

\fieldnotes{Warmup, then cosine to a tenth of the peak, linear decay to zero, or a constant peak with a cooldown over the last $10$--$20\%$ (WSD).}
{Convex Lipschitz SGD; the river valley, a slow river between stiff noisy hills.}
{A last-iterate bound for every schedule; linear decay attains the averaging rate $DG/\sqrt T$ exactly; exact second moments and the annealing decomposition in the river valley.}
{Loss $=$ area term $+$ recent-schedule term; a noise floor $\approx\eta\tr\Sigma_{\mathrm{hill}}/(4B)$ that a cooldown removes, leaving $\propto(\eta hK)^{-1/2}$; best cooldown fraction $\propto T^{(2\alpha-1)/3}$; in the bound, a floor $\rho$ costs $\rho\ln T$.}
{A landscape that moves with the schedule; position-dependent noise; Adam, momentum, weight decay; tuned peaks; why $20\%$.}
{Decay to (nearly) zero (proved: linear decay is worst-case optimal, and in the river valley a final step $\rho\eta$ keeps about a fraction $\rho$ of the floor); expect the cooldown drop to scale with the peak and $1/B$, and cool down for longer than the slowest relaxation time (derived); the exact $20\%$ is folklore backed by experiments.}

The two idealizations agree on three points. The learning rate should end at or near zero: the bound charges a final floor $\rho$ a term $\rho\ln T$, and the river valley leaves a fraction $\rho$ of the noise floor in place. The loss at the end is an area term minus a recent-schedule term, which is why very different shapes end close together and why the empirical laws work. And the benefit of a longer cooldown saturates. They disagree on its length: the bound, with a tuned peak, wants a decay over the whole run; the river valley wants a cooldown set by the relaxation times and the river's slope, a fixed fraction of the run only if the river falls like the inverse square root of the area. A short cooldown branched from a stable checkpoint measures the noise floor, the size of the prize, before the shape is tuned. That $20\%$ is enough is supported by experiments \parencite{HageleEtAl2024}, not by these theories.

The chapter closes a thread from \cref{ch:prologue}: sign descent never converges at a fixed step size, and only a schedule removes its two-cycle (Exercise~\ref{ex:schedules:sign}). Momentum and schedule-free methods remove such floors by averaging instead (\cref{ch:momentum}); the river valley is the noisy quadratic model of \cref{ch:noise} at the time scale of a run; and with decoupled weight decay the schedule is the step size of a Frank--Wolfe method, a cooldown its decaying step (\cref{ch:decay}).

\begin{exercises}
\item \consolidation (Two exact bounds.) Prove \cref{cor:schedules:shapes}, and show that at $T=10^4$ the constant step $D/(G\sqrt T)$ gives a bound $1.79$ times that of the best constant step.\label{ex:schedules:exact}

\item \consolidation (Shapes of the cooldown.) For $\varphi=1$ on $[0,1-c]$ and $\varphi(s)=((1-s)/c)^p$ after, show that $\Phi(0)=1-cp/(p+1)$ and $J[\varphi]=1+1/p+\ln((p+1-cp)/c)$. For $c=0.2$, find the $p$ that minimizes $\sqrt{J/\Phi(0)}$ (about $1.35$) and compare with the $1-\sqrt\cdot$ cooldown of \cref{tab:schedules:bound}. \hint{with $u=(1-s)/c$, the integrand of \eqref{eq:schedules:J} on the cooldown is $(p+1)^2u^{p-1}/(c(2p+1))$.}\label{ex:schedules:shape}

\item \consolidation (A floor costs a logarithm.) Let $\eta_t=\eta\varphi(t/(T+1))$ with $\varphi$ positive and Lipschitz on $[0,1]$ and $\varphi(1)=\rho$. Show that the last-iterate term of \eqref{eq:schedules:lastiterate} is $\frac12G^2\eta(\rho\ln T+O(1))$, and check the entry $1.775+0.1\ln T$ of \cref{tab:schedules:bound}. \hint{near the end, $R_k\approx\rho\eta(T-k+1)$ and $V_k\approx\rho^2\eta^2(T-k+1)$.}\label{ex:schedules:floor}

\item \consolidation (Step decay.) A hill starts at $v_i(\eta)$ and the step size drops to $\eta/10$. (a) Show that the excess over the new floor is below $1\%$ of the old floor after about $22.5/(\eta h)$ steps. (b) Show that an increase of the step size enters \eqref{eq:schedules:annealing} with a positive sign; describe the hill during a linear warmup.\label{ex:schedules:step}

\item \consolidation (The cooldown law.) (a) Derive \eqref{eq:schedules:cooldownlaw} from $dv/dA=-2hv+\eta(A)\sigma^2$. (b) Show that $\chi(\kappa)=1-\kappa/3+O(\kappa^2)$ and $\chi(\kappa)\sim\frac{\sqrt\pi}2\kappa^{-1/2}$. (c) Derive \eqref{eq:schedules:Kstar}, and show that at $K^\star$ the residual noise is twice the area cost $\omega\eta K^\star/2$.\label{ex:schedules:law}

\item \variation (Weighted conversions; \textcite{DefazioEtAl2023}.) Let $w_t>0$, $W=\sum_{t\le T}w_t$, $W_t=\sum_{s=t}^Tw_s$, $z_{t+1}=z_t-\gamma w_tg_t$ and $\eta_t=\gamma w_tW_{t+1}/W$. (a) Show that $W\theta_T=\sum_tw_tz_t$, and generalize the proof of \cref{prop:schedules:linear} to $\E\Loss(\theta_T)-\Loss^\star\le W^{-1}\bigl(D^2/(2\gamma)+\frac\gamma2G^2\sum_tw_t^2\bigr)$. (b) Show by the Cauchy--Schwarz inequality that uniform weights, that is linear decay, are best. (c) If $\E\norm{g_t}_2^2\le G_t^2$, show that the best weights are $w_t\propto G_t^{-2}$; what does $\eta_t\propto G_t^{-2}\sum_{s>t}G_s^{-2}$ look like when gradient norms fall?\label{ex:schedules:weights}

\item \variation (Normalized steps need a schedule.) Sign descent on $\frac12hx^2$ is $x_{t+1}=x_t-\eta_t\sign(x_t)$ (\cref{ch:prologue}). (a) With a constant step, show that the iterates end on a two-cycle with $\abs x\le\eta$, or at the fixed point $0$ when $x_1/\eta$ is an integer. (b) Let $\eta_t\ge\eta_{t-1}/2$ for all $t$; show that $\abs{x_t}\le\eta_{t-1}$ implies $\abs{x_{t+1}}\le\eta_t$, so the loss then stays below $h\eta_t^2/2$. (c) Compare this floor, quadratic in the step and present without noise, with the noise floor of \cref{cor:schedules:floor}.\label{ex:schedules:sign}

\item \variation (The noise-optimal cooldown.) In one hill, choose $\eta_t\ge0$ to minimize $v_K$ in $v_{t+1}=(1-\eta_th)^2v_t+\eta_t^2\sigma^2$. (a) Show that the step minimizing $v_{t+1}$ given $v_t$ is $hv_t/(\sigma^2+h^2v_t)$ and gives $1/v_{t+1}=1/v_t+h^2/\sigma^2$. (b) Show that this greedy choice is optimal. \hint{$v\mapsto\min_\eta\{(1-\eta h)^2v+\eta^2\sigma^2\}$ is increasing.} (c) Deduce $v_K=\sigma^2/(h^2(K+n_0))$ with $n_0=\sigma^2/(h^2v_0)$, reached by $\eta_t=1/(h(t+n_0+1))$, and compare with \eqref{eq:schedules:cooldownlaw}.\label{ex:schedules:greedy}

\item \variation (Averaging instead of cooling.) In one hill at equilibrium at the step $\eta$, let $\bar y$ average the last $K$ iterates; put $b=1-\eta h$ and $v=v_i(\eta)$. (a) Show that $\E\bar y^2=\frac vK\bigl(\frac{1+b}{1-b}-\frac{2b(1-b^K)}{K(1-b)^2}\bigr)$, so that the hill loss at $\bar y$ is about $2/(\eta hK)$ times the floor. (b) Show that with a linear river the average gives up the same area as a linear cooldown of $K$ steps, up to one step. (c) Conclude that tail averaging beats a linear cooldown when $\eta hK$ is large; compare with \cref{ch:momentum} and with \textcite{HageleEtAl2024}, who compare averaging with cooldowns.\label{ex:schedules:averaging}

\item \simulation (Shapes and peaks.) In \texttt{12-schedules-river.py}: (a) replace the linear cooldown by a $1-\sqrt\cdot$ cooldown of the same length and check that the final loss falls from $0.131$ to $0.129$; (b) raise the peak to $0.2$, check that cosine ($0.080$) now beats WSD ($0.096$), and explain why with the river and hill parts; (c) add a linear warmup over $100$ steps and check that, at the end of the run, only the river parts change.\label{ex:schedules:sim1}

\item \simulation (The cone.) In \texttt{12-schedules-bound.py}: (a) check that constant steps settle at $\eta G^2/2$, and that after the step is halved the excess decays by a factor $e$ in about $130$ steps, at $\eta=0.0032$ and at $\eta=0.01$; (b) set $d=400$ and explain, with the first term of \eqref{eq:schedules:lastiterate}, why cosine and linear decay now end above constant steps.\label{ex:schedules:sim2}

\item \research (Derive the river.) Find a network model on structured data in which a river valley, with its exponent $\alpha$ and hill spectrum, can be derived from the data. \Textcite{WenEtAl2025} give a bigram example; a model with learned features is open.\label{ex:schedules:derive}

\item \research (Steps or area?) The river valley predicts that the loss relaxes after a decrease of the learning rate in area, the cone that it relaxes in steps. Which mechanism (a nonsmooth minimum, curvature adapting at the edge of stability, Adam's normalization) produces the learning-rate dependence that the empirical laws fit, and with what exponent? This is open, as far as we know.\label{ex:schedules:relax}
\end{exercises}

\notesandrefs
The step-size conditions of \textcite{RobbinsMonro1951} and the averaging of \textcite{PolyakJuditsky1992} are the classical theory of schedules, and \textcite{NemirovskiEtAl2009} showed that constant steps with averaging are optimal for convex Lipschitz problems. For the last iterate, \textcite{ShamirZhang2013} proved the $\log T/\sqrt T$ rate of the steps $1/\sqrt t$, \textcite{HarveyEtAl2019} showed it tight, \textcite{JainNagarajNetrapalli2019} designed a horizon-dependent schedule with the optimal rate, \textcite{Orabona2020} removed the bounded domain, and \textcite{LiuZhou2024} extended the analysis further. \Textcite{ZamaniGlineur2025} found linear decay optimal for the subgradient method, and \textcite{DefazioEtAl2023} for stochastic gradient descent. Our proof of \cref{thm:schedules:lastiterate} follows the comparison argument; our proof of \cref{prop:schedules:linear} is a conversion of the kind used by \textcite{DefazioEtAl2023}, written out for uniform weights.

WSD was introduced for MiniCPM \parencite{HuEtAl2024} and studied by \textcite{HageleEtAl2024}, and \textcite{SchaippEtAl2025} found its agreement with the convex bound; the river valley is due to \textcite{WenEtAl2025}, the laws of loss curves across schedules to \textcite{TissueEtAl2025} and \textcite{LuoEtAl2025}. We have not found the cooldown law \eqref{eq:schedules:cooldownlaw}, the optimal cooldowns \eqref{eq:schedules:Kstar} and \eqref{eq:schedules:cstar}, or the gamma-type saturation in the literature we could check; all are derived here, not proved.

\paragraph{Translating the literature.} \Textcite{DefazioEtAl2023} and \textcite{SchaippEtAl2025} write $x_t$ for our $\theta_t$ and $f$ for $\Loss$, and start at $x_1$ as we do. \Textcite{SchaippEtAl2025} write the schedule as a base learning rate $\gamma$ times a shape $\eta_t$, our peak $\eta$ times $\varphi(t/T)$, call the bound $\Omega_t$, allow a bound $G_t^2\ge\E\norm{g_t}^2$ that varies with $t$, and write $D=\norm{x_1-x_\star}$; their Theorem~3.1, which they present as a version of Theorem~10 of \textcite{DefazioEtAl2023}, is built from the same sums as \eqref{eq:schedules:lastiterate}, a term in $D^2$ and the squared steps and a term with the steps still to come, with each $\eta_t^2$ weighted by its own $G_t^2$. \Textcite{DefazioEtAl2023} work with weights $w_t$; their schedule $\eta_t\propto w_t\sum_{s>t}w_s$ is that of Exercise~\ref{ex:schedules:weights}. In \textcite{TissueEtAl2025} and \textcite{LuoEtAl2025} the forward area $S_1(t)=\sum_{s\le t}\eta_s$ is our $A_{t+1}$, the area $\sum_{s=k}^{t}\eta_s$ since a decrease at step $k$ is our $A_{t+1}-A_k$, the loss drop $\mathrm{LD}(t)$ is our annealing sum in \eqref{eq:schedules:law}, their $L_0$ and $A$ are our $L_\infty$ and $C$, and their $\alpha$ is our river exponent. The decay factor $\lambda$ (typically $0.999$) in the annealing area of \textcite{TissueEtAl2025} is our $e^{-2\bar\eta h}$, and the exponent of the power-law saturation of \textcite{LuoEtAl2025} plays the role of our gamma shape $\nu$.

\printchapterbib
\end{refsection}
